%==============================================================
\section{李代数的分类}
%==============================================================

\subsection{单李代数与半单李代数}

\textbf{单李代数：} 没有非平凡理想（$[I, \mathfrak{g}] \subset I$ 的子空间 $I$）。

\textbf{半单李代数：} 单李代数的直和。

\textbf{分类定理（Cartan--Killing）：} 复半单李代数完全由 Dynkin 图分类。

\subsection{经典李代数系列}

\begin{center}
\begin{tabular}{p{2cm}p{3cm}p{2cm}p{5cm}}
\toprule
\textbf{类型} & \textbf{李代数} & \textbf{维度} & \textbf{物理} \\
\midrule
$A_n$ & $\mathfrak{sl}(n+1, \mathbb{C})$ & $n(n+2)$ & $SU(n+1)$：色、味对称 \\
$B_n$ & $\mathfrak{so}(2n+1, \mathbb{C})$ & $n(2n+1)$ & $SO(2n+1)$：旋转 \\
$C_n$ & $\mathfrak{sp}(2n, \mathbb{C})$ & $n(2n+1)$ & $Sp(n)$：辛对称 \\
$D_n$ & $\mathfrak{so}(2n, \mathbb{C})$ & $n(2n-1)$ & $SO(2n)$：旋转 \\
\bottomrule
\end{tabular}
\end{center}

\subsection{例外李代数}

\begin{center}
\begin{tabular}{p{2cm}p{2cm}p{6cm}}
\toprule
\textbf{类型} & \textbf{维度} & \textbf{出现} \\
\midrule
$G_2$ & 14 & 八元数的自同构群；弦论 \\
$F_4$ & 52 & 例外 Jordan 代数；弦论 \\
$E_6$ & 78 & 大统一理论候选 \\
$E_7$ & 133 & 超引力；弦论 \\
$E_8$ & 248 & 弦论（$E_8 \times E_8$ 杂化弦）；凝聚态（Ising 模型） \\
\bottomrule
\end{tabular}
\end{center}

\subsection{根系与 Dynkin 图}

\subsubsection{Cartan 子代数与根}

\textbf{Cartan 子代数 $\mathfrak{h}$：} 极大交换子代数（所有元素可同时对角化）。

\textbf{根：} $\alpha \in \mathfrak{h}^*$ 使得存在 $E_\alpha \neq 0$ 满足 $[H, E_\alpha] = \alpha(H)E_\alpha$，$\forall H \in \mathfrak{h}$。

\textbf{分解：}
\[
\mathfrak{g} = \mathfrak{h} \oplus \bigoplus_{\alpha \in \Phi} \mathfrak{g}_\alpha
\]

\subsubsection{Dynkin 图}

\textbf{构造规则：}
\begin{itemize}
  \item 每个单根画一个节点
  \item 两节点间的边数由夹角决定：$0° \to 0$ 边，$60° \to 1$ 边，$45° \to 2$ 边，$30° \to 3$ 边
  \item 不等长根之间加箭头（指向短根）
\end{itemize}

\textbf{经典系列的 Dynkin 图：}

\begin{itemize}
  \item $A_n$：$\circ - \circ - \circ - \cdots - \circ$（$n$ 个节点，链状）
  \item $B_n$：$\circ - \circ - \cdots - \circ \Rightarrow \circ$（末端双键）
  \item $C_n$：$\circ - \circ - \cdots - \circ \Leftarrow \circ$（末端双键，方向反）
  \item $D_n$：$\circ - \circ - \cdots - \circ$ 末端分叉为两个节点
\end{itemize}

\textbf{例外系列：} $G_2$、$F_4$、$E_6$、$E_7$、$E_8$ 有特殊的 Dynkin 图。

\subsection{表示论}

\subsubsection{权与最高权}

\textbf{权：} 表示空间中 $H \in \mathfrak{h}$ 的特征值。

\textbf{最高权：} 在某个偏序下最大的权。

\textbf{定理：} 不可约表示由最高权唯一确定。

\subsubsection{$\mathfrak{su}(2)$ 的表示}

不可约表示由非负整数/半整数 $j$ 标记，维度 $2j + 1$：
\begin{itemize}
  \item $j = 0$：标量（1 维）
  \item $j = 1/2$：旋量（2 维）——电子自旋
  \item $j = 1$：向量（3 维）——光子
  \item $j = 3/2$：（4 维）——$\Delta$ 重子
  \item $j = 2$：（5 维）——引力子
\end{itemize}

\subsubsection{$\mathfrak{su}(3)$ 的表示与粒子物理}

\textbf{基本表示 $\mathbf{3}$：} 夸克（$u, d, s$）

\textbf{共轭表示 $\bar{\mathbf{3}}$：} 反夸克

\textbf{分解：}
\[
\mathbf{3} \otimes \bar{\mathbf{3}} = \mathbf{8} \oplus \mathbf{1}
\]

$\mathbf{8}$：8 个胶子（伴随表示）；$\mathbf{1}$：单态。

\[
\mathbf{3} \otimes \mathbf{3} \otimes \mathbf{3} = \mathbf{10} \oplus \mathbf{8} \oplus \mathbf{8} \oplus \mathbf{1}
\]

$\mathbf{10}$：重子十重态（$\Delta$、$\Sigma^*$、$\Xi^*$、$\Omega^-$）——这是盖尔曼预言 $\Omega^-$ 的基础。

%==============================================================
\section{拓扑结构}
%==============================================================

\subsection{基本群}

\textbf{基本群 $\pi_1(X)$：} 闭合路径的同伦类。

\begin{center}
\begin{tabular}{p{3cm}p{3cm}p{5cm}}
\toprule
\textbf{空间} & \textbf{$\pi_1$} & \textbf{物理意义} \\
\midrule
$S^1$ & $\mathbb{Z}$ & 绕数（Aharonov--Bohm 效应） \\
$SO(3) \cong \mathbb{RP}^3$ & $\mathbb{Z}_2$ & 自旋 $1/2$（$2\pi$ 旋转 $\neq$ 恒等） \\
$SU(2) \cong S^3$ & $0$（单连通） & 万有覆盖 \\
$U(1) \cong S^1$ & $\mathbb{Z}$ & 磁通量子化、涡旋 \\
$SO(3,1)^+$（Lorentz） & $\mathbb{Z}_2$ & 旋量 \\
\bottomrule
\end{tabular}
\end{center}

\subsection{高阶同伦群}

\textbf{$\pi_n(X)$：} $S^n \to X$ 的连续映射的同伦类。

\begin{center}
\begin{tabular}{p{3cm}p{3cm}p{5cm}}
\toprule
\textbf{同伦群} & \textbf{值} & \textbf{物理意义} \\
\midrule
$\pi_2(S^2)$ & $\mathbb{Z}$ & 磁单极（'t Hooft--Polyakov） \\
$\pi_3(S^2)$ & $\mathbb{Z}$ & Skyrmion（核子模型） \\
$\pi_3(SU(2))$ & $\mathbb{Z}$ & 瞬子（Yang--Mills） \\
$\pi_3(SU(3))$ & $\mathbb{Z}$ & QCD 真空结构（$\theta$ 真空） \\
$\pi_1(U(1))$ & $\mathbb{Z}$ & 涡旋（超流、超导） \\
\bottomrule
\end{tabular}
\end{center}

\subsection{纤维丛与拓扑}

\textbf{主 $G$-丛：} $P \xrightarrow{\pi} M$，纤维是群 $G$。

\textbf{分类：} 主 $G$-丛的同构类由 $[M, BG]$ 给出（$BG$ 是分类空间）。

\textbf{特征类：}
\begin{itemize}
  \item $U(1)$ 丛：第一 Chern 类 $c_1 \in H^2(M, \mathbb{Z})$ → 磁通量子化
  \item $SU(n)$ 丛：Chern 类 $c_k$ → 瞬子数
  \item $SO(n)$ 丛：Pontryagin 类 $p_k$、Euler 类 $e$
\end{itemize}

\textbf{在物理中：}
\begin{itemize}
  \item 电磁场是 $U(1)$ 主丛上的联络
  \item Yang--Mills 场是 $SU(n)$ 主丛上的联络
  \item 瞬子数 = 第二 Chern 数 = $\frac{1}{8\pi^2}\int \text{tr}(F \wedge F)$
\end{itemize}

\subsection{de Rham 上同调}

\textbf{de Rham 上同调：}
\[
H^k_{dR}(M) = \frac{\ker(d: \Omega^k \to \Omega^{k+1})}{\text{im}(d: \Omega^{k-1} \to \Omega^k)} = \frac{\text{闭形式}}{\text{恰当形式}}
\]

\textbf{物理意义：}
\begin{itemize}
  \item $H^2_{dR}(M) \neq 0$：存在非平凡的电磁场构型（磁单极）
  \item $H^1_{dR}(M) \neq 0$：存在非平凡的矢势（Aharonov--Bohm）
\end{itemize}

\subsection{同调与流形}

\textbf{Euler 示性数：}
\[
\chi(M) = \sum_{k=0}^n (-1)^k b_k
\]

其中 $b_k = \dim H^k(M)$ 是 Betti 数。

\textbf{Gauss--Bonnet 定理：}
\[
\int_M K\,dA = 2\pi\chi(M)
\]

（$K$ 是高斯曲率，$M$ 是闭曲面）

\textbf{在 CG 中：} 网格的拓扑不变量（亏格）影响参数化和网格生成。